\section{反向 SDE 与生成过程}

\subsection{Anderson 定理：反向 SDE}

给定前向 SDE

$$
\mathrm{d}\mathbf{x} = f(\mathbf{x}, t)\,\mathrm{d}t + g(t)\,\mathrm{d}\mathbf{w}
$$

Anderson（1982）证明了对应的反向时间 SDE 为：

$$
\mathrm{d}\mathbf{x} = \left[f(\mathbf{x}, t) - g(t)^2 \nabla_{\mathbf{x}} \log p_t(\mathbf{x})\right]\mathrm{d}t + g(t)\,\mathrm{d}\bar{\mathbf{w}}
$$

其中 $\bar{\mathbf{w}}$ 是反向时间的 Wiener 过程，$\mathrm{d}t$ 在此为\textbf{负}的时间增量（时间从 $T$ 走回 $0$）。

\begin{importantbox}{反向 SDE 的结构}
反向 SDE 与前向 SDE 有着惊人的相似结构：
\begin{itemize}
    \item 扩散系数 $g(t)$ \textbf{完全相同}；
    \item 漂移项在原有 $f(\mathbf{x}, t)$ 基础上\textbf{额外减去} $g(t)^2 \nabla_{\mathbf{x}} \log p_t(\mathbf{x})$。
\end{itemize}
唯一需要额外知道的就是 \textbf{score function} $\nabla_{\mathbf{x}} \log p_t(\mathbf{x})$——这正是扩散模型通过神经网络学习的目标。
\end{importantbox}

\subsection{Score Function 的角色}

回顾离散 DDPM 中的做法：
\begin{enumerate}
    \item 定义前向过程（加噪）；
    \item 训练噪声预测网络 $\boldsymbol{\epsilon}_\theta(\mathbf{x}_t, t)$；
    \item 通过 $\nabla_{\mathbf{x}} \log p_t(\mathbf{x}_t) = -\frac{\boldsymbol{\epsilon}}{\sigma_t}$ 的关系获得 score；
    \item 利用 score 执行反向去噪过程。
\end{enumerate}

在连续时间框架下，逻辑完全一致：
\begin{enumerate}
    \item 定义前向 SDE（漂移 + 扩散）；
    \item 训练 score 网络 $\mathbf{s}_\theta(\mathbf{x}, t) \approx \nabla_{\mathbf{x}} \log p_t(\mathbf{x})$；
    \item 将 $\mathbf{s}_\theta$ 代入反向 SDE，即可执行生成。
\end{enumerate}

\begin{knowledgebox}{Score 与噪声预测的等价性}
对于跳跃分布 $q(\mathbf{x}_t | \mathbf{x}_0) = \mathcal{N}(\alpha_t \mathbf{x}_0, \sigma_t^2 \mathbf{I})$，score function 可以表示为：
$$
\nabla_{\mathbf{x}_t} \log p_t(\mathbf{x}_t | \mathbf{x}_0) = -\frac{\mathbf{x}_t - \alpha_t \mathbf{x}_0}{\sigma_t^2} = -\frac{\boldsymbol{\epsilon}}{\sigma_t}
$$
其中 $\boldsymbol{\epsilon}$ 是加噪时采样的标准高斯噪声。因此，预测 score 等价于预测噪声（除以 $\sigma_t$ 的差别）。这就是为什么 DDPM 的"噪声预测"训练目标和 Score Matching 的训练目标本质上是一回事。
\end{knowledgebox}

\subsection{反向过程的离散化实现}

虽然反向 SDE 是连续的，实际计算中需要离散化。给定反向 SDE，我们将时间 $[0, T]$ 划分为 $N$ 个步骤，在每一步中：

$$
\mathbf{x}_{t-\Delta t} \approx \mathbf{x}_t + \left[f(\mathbf{x}_t, t) - g(t)^2 \mathbf{s}_\theta(\mathbf{x}_t, t)\right](-\Delta t) + g(t)\sqrt{\Delta t}\;\boldsymbol{\epsilon}
$$

这与 DDPM 的离散反向过程完全一致：每一步使用 score（或噪声预测）修正均值，再添加适当的随机噪声。

\begin{warningbox}{离散化步数与生成质量}
SDE 的离散化本质上是 Euler-Maruyama 方法。步数越多，近似越精确，生成质量越好——但计算成本也越高。DDPM 需要约 1000 步才能获得良好质量，而 DDIM（对应确定性 ODE 的离散化）只需约 50 步。后续讲座将介绍更高效的 SDE/ODE solver（如 DPM-Solver），可以将步数进一步减少到 10--15 步。
\end{warningbox}

\subsection{本章小结}

Anderson 定理保证了：给定任意前向 SDE，反向 SDE 自动确定——唯一需要学习的就是 score function。离散化反向 SDE 后，我们恢复了经典扩散模型的采样算法。


